% =============================================================================
% Individual-level DiD results: main College x Post DiD, heterogeneity, and
% placebo counties. Conway, Iselin & Rork, "Mult-No-More?"
%
% Compile from the Overleaf project root (next to figures/ and tables/, the
% folders mirrored from results/overleaf_export/). Every number in the text
% comes from the in-text numbers register via \pn{<key>}; unknown keys print ??.
%
% Slide titles are placeholders: [TITLE: what the slide shows].
% Outputs used (all from code/stata/02_did_analysis.do, 02_did_placebo.do):
%   tables/tab_did_main.tex, tables/tab_did_heterogeneity.tex,
%   tables/tab_did_placebo.tex, figures/fig_es_*_migration.png,
%   figures/fig_es_age_*_migration.png, figures/fig_did_placebo.pdf
% =============================================================================
\documentclass[aspectratio=169,10pt]{beamer}
\usetheme{default}
\setbeamertemplate{navigation symbols}{}
\setbeamertemplate{footline}[frame number]
\usepackage{booktabs}
\usepackage{graphicx}
\usepackage{amsmath}
\input{tables/paper_numbers}

\newcommand{\slidenote}[1]{\par\vfill{\scriptsize #1}}

\title{[TITLE: deck title]}
\author{Conway, Iselin \& Rork}
\date{}

\begin{document}

\begin{frame}
  \titlepage
\end{frame}

% -----------------------------------------------------------------------------
\begin{frame}{[TITLE: DiD design]}
  \begin{itemize}
    \item ACS person-level microdata, 2016--2024; adults 25+; 2020 excluded
    \item Treatment proxy: \textbf{College degree $\times$ Post-2020}
      \begin{itemize}
        \item Preschool for All's thresholds (\$125k single / \$200k joint) fall
              mostly on college graduates
      \end{itemize}
    \item Outcomes (percentage points)
      \begin{itemize}
        \item Out-migration and out-of-state moves: last year's Multnomah residents
        \item In-migration: residents of the West Coast, or of the Lower 48 + DC,
              who moved into Multnomah
      \end{itemize}
    \item Fixed effects: year and demographic cells (age, sex, marital status,
          children, education); origin county for in-migration
    \item Robust SEs (out-migration); clustered by origin county (in-migration)
  \end{itemize}
\end{frame}

% -----------------------------------------------------------------------------
\begin{frame}{[TITLE: main DiD table]}
  \centering
  \resizebox{0.9\textwidth}{!}{\input{tables/tab_did_main}}
  \slidenote{College $\times$ Post coefficient, percentage points. $^{*}p<0.10$,
  $^{**}p<0.05$, $^{***}p<0.01$. Implied semi-elasticities:
  out-migration \pn{did_semi_elast_out_county}, out-of-state \pn{did_semi_elast_out_state}.}
\end{frame}

% -----------------------------------------------------------------------------
\begin{frame}{[TITLE: DiD event studies]}
  \centering
  \begin{tabular}{@{}cc@{}}
    \scriptsize Out-migration & \scriptsize Out-of-state \\
    \includegraphics[height=0.33\textheight]{figures/fig_es_out_migration.png} &
    \includegraphics[height=0.33\textheight]{figures/fig_es_out_state_migration.png} \\
    \scriptsize In-migration (West Coast) & \scriptsize In-migration (Lower 48 + DC) \\
    \includegraphics[height=0.33\textheight]{figures/fig_es_in_migration_west.png} &
    \includegraphics[height=0.33\textheight]{figures/fig_es_in_migration_48.png}
  \end{tabular}
  \slidenote{College $\times$ Year coefficients relative to 2019 (2020 excluded), 95\% CIs.
  Post-2020 out-migration coefficients range from \pn{did_es_out_college_post_min}
  to \pn{did_es_out_college_post_max} pp.}
\end{frame}

% -----------------------------------------------------------------------------
\begin{frame}{[TITLE: heterogeneity table]}
  \centering
  \resizebox*{!}{0.78\textheight}{\input{tables/tab_did_heterogeneity}}
  \slidenote{College $\times$ Post $\times$ group triple differences, with year $\times$ group and
  college $\times$ group fixed effects. WFH: Dingel--Neiman teleworkability of the
  respondent's occupation; ``no occupation'' is mostly retirees and non-workers.}
\end{frame}

% -----------------------------------------------------------------------------
\begin{frame}{[TITLE: heterogeneity takeaways]}
  \begin{itemize}
    \item Age: out-migration response concentrated among 45--64 year-olds
    \item Larger for those without own children and for the unmarried
          (differences not significant)
    \item WFH: response \textbf{larger in non-teleworkable occupations}; near zero
          with no occupation
      \begin{itemize}
        \item Not consistent with a ``remote workers left'' story
      \end{itemize}
    \item No heterogeneity in in-migration: effects near zero for every group
  \end{itemize}
  \slidenote{Each statement reads off the heterogeneity table; equal-effects $p$-values are in
  its last row per panel.}
\end{frame}

% -----------------------------------------------------------------------------
\begin{frame}{[TITLE: event studies by age]}
  \centering
  \begin{tabular}{@{}cc@{}}
    \scriptsize Out-migration & \scriptsize Out-of-state \\
    \includegraphics[height=0.6\textheight]{figures/fig_es_age_out_migration.png} &
    \includegraphics[height=0.6\textheight]{figures/fig_es_age_out_state_migration.png}
  \end{tabular}
  \slidenote{Ages 45--64: post-2022 out-migration coefficients between
  \pn{did_es_out_college_age45_64_post22_min} and \pn{did_es_out_college_age45_64_post22_max} pp.}
\end{frame}

% -----------------------------------------------------------------------------
\begin{frame}{[TITLE: event studies by own children]}
  \centering
  \begin{tabular}{@{}cc@{}}
    \scriptsize Out-migration & \scriptsize Out-of-state \\
    \includegraphics[height=0.6\textheight]{figures/fig_es_kids_out_migration.png} &
    \includegraphics[height=0.6\textheight]{figures/fig_es_kids_out_state_migration.png}
  \end{tabular}
  \slidenote{College $\times$ Year $\times$ group coefficients relative to 2019, 95\% CIs. Legend: each group's
  pre-trend $p$-value (2016--18 coefficients jointly zero).}
\end{frame}

% -----------------------------------------------------------------------------
\begin{frame}{[TITLE: event studies by marital status]}
  \centering
  \begin{tabular}{@{}cc@{}}
    \scriptsize Out-migration & \scriptsize Out-of-state \\
    \includegraphics[height=0.6\textheight]{figures/fig_es_married_out_migration.png} &
    \includegraphics[height=0.6\textheight]{figures/fig_es_married_out_state_migration.png}
  \end{tabular}
  \slidenote{College $\times$ Year $\times$ group coefficients relative to 2019, 95\% CIs. Legend: each group's
  pre-trend $p$-value (2016--18 coefficients jointly zero).}
\end{frame}

% -----------------------------------------------------------------------------
\begin{frame}{[TITLE: event studies by WFH potential]}
  \centering
  \begin{tabular}{@{}cc@{}}
    \scriptsize Out-migration & \scriptsize Out-of-state \\
    \includegraphics[height=0.6\textheight]{figures/fig_es_wfh_out_migration.png} &
    \includegraphics[height=0.6\textheight]{figures/fig_es_wfh_out_state_migration.png}
  \end{tabular}
  \slidenote{College $\times$ Year $\times$ group coefficients relative to 2019, 95\% CIs. Legend: each group's
  pre-trend $p$-value (2016--18 coefficients jointly zero).}
\end{frame}

% -----------------------------------------------------------------------------
\begin{frame}{[TITLE: placebo design]}
  \begin{itemize}
    \item Rerun the main DiD as if each Policy-match donor county were treated
      \begin{itemize}
        \item \pn{did_placebo_n_candidates} donors identifiable as ACS origin counties
      \end{itemize}
    \item Keep counties with reasonable pre-2020 trends --- two screens:
      \begin{enumerate}
        \item \textbf{Within county:} its own 2016--18 College $\times$ Year
              coefficients jointly insignificant ($p > 0.10$)
        \item \textbf{Like Multnomah:} 2016--19 IRS AGI out- and in-migration rate
              slopes among the closest half of donors to Multnomah's
      \end{enumerate}
    \item \pn{did_placebo_n_screened} counties pass both
    \item Inference: one-sided permutation $p$-value, the share of placebos (plus
          Multnomah) with a response at least as large as Multnomah's
  \end{itemize}
\end{frame}

% -----------------------------------------------------------------------------
\begin{frame}{[TITLE: placebo estimates]}
  \centering
  \includegraphics[width=\textwidth]{figures/fig_did_placebo.pdf}
  \slidenote{One-sided permutation $p$-values, screened placebos: out-migration
  \pn{did_placebo_perm_p_out_screened}, out-of-state \pn{did_placebo_perm_p_out_state_screened},
  in-migration \pn{did_placebo_perm_p_in_48_screened}. All \pn{did_placebo_n_candidates} donors:
  \pn{did_placebo_perm_p_out_all}, \pn{did_placebo_perm_p_out_state_all},
  \pn{did_placebo_perm_p_in_48_all}. With \pn{did_placebo_n_screened} placebos the smallest
  attainable $p$-value is $1/(\text{placebos}+1)$.}
\end{frame}

% -----------------------------------------------------------------------------
\begin{frame}{[TITLE: placebo event studies]}
  \centering
  \includegraphics[width=\textwidth]{figures/fig_did_placebo_es.pdf}
  \slidenote{College $\times$ Year coefficients relative to 2019: Multnomah (with 95\% CIs) and the
  \pn{did_placebo_n_screened} screened placebo counties. 2020 excluded.}
\end{frame}

% -----------------------------------------------------------------------------
\begin{frame}{[TITLE: Orleans Parish]}
  \begin{itemize}
    \item Orleans Parish, LA passes both screens but shows a large
          \textbf{negative} College $\times$ Post response (see placebo table)
    \item {[DISCUSS: why college out-migration from New Orleans fell after 2020]}
    \item A two-sided test would count Orleans against Multnomah; the one-sided
          test asks only whether a similar county saw college out-migration rise as much
  \end{itemize}
\end{frame}

% -----------------------------------------------------------------------------
\appendix
\begin{frame}{[TITLE: placebo table (backup)]}
  \centering
  \resizebox*{!}{0.8\textheight}{\input{tables/tab_did_placebo}}
  \slidenote{New Jersey raised its rate on income of \$1--5M from 8.97\% to 10.75\% from tax year
  2020; New York added 9.65--10.9\% brackets above about \$1.1M (single) from tax year 2021.
  Excluding NJ and NY counties, one-sided $p$-values (screened):
  \pn{did_placebo_perm_p_out_nonjny}, \pn{did_placebo_perm_p_out_state_nonjny},
  \pn{did_placebo_perm_p_in_48_nonjny}.}
\end{frame}

\end{document}
